\section{Materials and methods}
\label{sec:methods}

\subsection{Culture conditions}

\emph{Thalassiosira pseudonana} strain \strain{} was grown in artificial
seawater at a salinity of 32 with f/2 trace metals and vitamins, at
$18.0 \pm 0.3\,^{\circ}\mathrm{C}$ under a 14:10 light and dark cycle at
$120\;\mu\mathrm{mol\,photons\,m^{-2}\,s^{-1}}$. The nitrogen-replete regime
supplied $880\;\mu\mathrm{mol\,L^{-1}}$ nitrate. The nitrogen-limited regime
supplied $55\;\mu\mathrm{mol\,L^{-1}}$ nitrate and was otherwise identical.
Three biological replicates per regime were carried in separate 2 L
borosilicate flasks with filtered air bubbling at
$40\;\mathrm{mL\,min^{-1}}$. Cultures were acclimated for six generations
before the experimental transfer.

Specific growth rate over the exponential interval was taken from cell counts
$X$ at times $t_1$ and $t_2$,
\begin{equation}
  \label{eq:growth}
  \mu = \frac{\ln X(t_2) - \ln X(t_1)}{t_2 - t_1},
\end{equation}
with counts made in triplicate on a haemocytometer at 24 hour intervals.
Table~\ref{tab:physiology} reports the resulting physiology.

\subsection{Sampling and RNA extraction protocol}

Sampling took place four hours into the light period on day 5, at
mid-exponential phase in both regimes. The protocol below was followed without
deviation for all six flasks.

\begin{enumerate}
  \item Draw 250 mL of culture into a pre-chilled flask and record the time.
  \item Filter onto a 47 mm, $0.8\;\mu\mathrm{m}$ polycarbonate membrane at a
        vacuum no stronger than 15 kPa, keeping total filtration time under
        four minutes.
  \item Transfer the membrane to a 2 mL screw-cap tube, flash freeze in liquid
        nitrogen, and hold at $-80\,^{\circ}\mathrm{C}$ for no longer than
        eleven days.
  \item Add 1 mL of guanidinium thiocyanate lysis buffer and 200 mg of
        acid-washed 0.5 mm glass beads. Bead beat for
        $3 \times 45\;\mathrm{s}$ at $6.0\;\mathrm{m\,s^{-1}}$ with 60 s on ice
        between cycles.
  \item Clear the lysate at $12{,}000 \times g$ for 5 minutes at
        $4\,^{\circ}\mathrm{C}$ and carry the supernatant onto a silica column.
  \item Perform on-column DNase I digestion for 20 minutes at room temperature,
        then wash twice and elute in $40\;\mu\mathrm{L}$ of nuclease-free water.
  \item Assess integrity on a microfluidic electrophoresis chip. Accept the
        sample only if the RNA integrity number is at least 8.0 and the
        $28\mathrm{S}/18\mathrm{S}$ ratio lies between 1.4 and 2.2.
\end{enumerate}

All six extractions passed step 7 on the first attempt, with integrity numbers
between 8.4 and 9.1 and yields between 3.1 and $7.8\;\mu\mathrm{g}$.

\subsection{Library preparation and sequencing}

Polyadenylated RNA was selected with oligo-dT magnetic beads from
$1.0\;\mu\mathrm{g}$ of total RNA, fragmented for 6 minutes at
$94\,^{\circ}\mathrm{C}$, and converted to stranded libraries by the dUTP
second-strand method. Libraries were amplified for 11 cycles, quantified by
fluorometry, pooled equimolar, and sequenced as $2 \times 150$ paired-end reads
on a single patterned flow cell lane. Two libraries, NL-2 and NL-3, were loaded
from a pool prepared one day later than the rest, which is the most likely
origin of their elevated duplication.

\begin{table}[t]
  \centering
  \caption{Physiological state at the sampling point. Values are the mean of
  three biological replicates with the standard deviation. Residual nitrate is
  the concentration remaining in the medium at sampling.}
  \label{tab:physiology}
  \begin{tabular}{lcc}
    \toprule
    Measurement & Nitrogen replete & Nitrogen limited \\
    \midrule
    Growth rate $\mu$ ($\mathrm{d^{-1}}$)
      & $1.34 \pm 0.05$ & $0.41 \pm 0.07$ \\
    Cell density ($10^{6}\;\mathrm{cells\,mL^{-1}}$)
      & $4.82 \pm 0.31$ & $1.96 \pm 0.24$ \\
    Chlorophyll $a$ ($\mathrm{pg\,cell^{-1}}$)
      & $0.191 \pm 0.012$ & $0.074 \pm 0.009$ \\
    Residual nitrate ($\mu\mathrm{mol\,L^{-1}}$)
      & $612 \pm 28$ & $1.8 \pm 0.6$ \\
    Particulate C:N (mol mol$^{-1}$)
      & $6.9 \pm 0.2$ & $14.7 \pm 1.1$ \\
    Maximum quantum yield $F_v/F_m$
      & $0.66 \pm 0.01$ & $0.38 \pm 0.03$ \\
    \bottomrule
  \end{tabular}
\end{table}

\subsection{Analysis pipeline}

Figure~\ref{fig:workflow} sets out the full workflow from flask to gene list.
The computational half was run once, end to end, on a fixed container image so
that the version of every tool is recoverable from the report. Adapters and low
quality tails were removed with fastp \citep{chen2018fastp}, reads were aligned
to reference assembly v3 with a splice-aware aligner \citep{dobin2013star},
fragments were assigned to transcripts with featureCounts
\citep{liao2014featurecounts}, and the per-stage metrics were collected into a
single report \citep{ewels2016multiqc}. Parameter choices follow the community
recommendations for a small genome with short introns
\citep{conesa2016survey}.

\begin{figure}[t]
  \centering
  \resizebox{\linewidth}{!}{\begin{tikzpicture}[
      font=\sffamily\scriptsize,
      wet/.style={draw=wetlab,line width=0.7pt,rounded corners=2pt,
        fill=wetlab!8,align=center,minimum width=20mm,minimum height=9mm,
        inner sep=3pt},
      dry/.style={draw=drylab,line width=0.7pt,rounded corners=2pt,
        fill=drylab!8,align=center,minimum width=20mm,minimum height=9mm,
        inner sep=3pt},
      gate/.style={draw=accent,line width=0.7pt,fill=white,align=center,
        minimum width=16mm,minimum height=8mm,inner sep=2pt},
      flow/.style={-{Stealth[length=2mm,width=1.6mm]},line width=0.6pt,
        draw=ink},
      lane/.style={font=\sffamily\footnotesize\bfseries},
    ]

    \node[wet] (culture) at (0,3.3) {two nitrogen\\ regimes};
    \node[wet] (sample) at (2.7,3.3) {filter and\\ flash freeze};
    \node[wet] (rna) at (5.4,3.3) {bead beat and\\ column extract};
    \node[gate] (rin) at (8.1,3.3) {RIN $\geq 8.0$};
    \node[wet] (library) at (10.8,3.3) {stranded\\ polyA library};
    \node[wet] (seq) at (13.5,3.3) {$2 \times 150$\\ sequencing};

    \node[dry] (trim) at (13.5,0.9) {adapter and\\ quality trim};
    \node[dry] (align) at (10.8,0.9) {splice-aware\\ alignment};
    \node[gate] (mapgate) at (8.1,0.9) {mapped\\ $\geq 85\%$};
    \node[dry] (count) at (5.4,0.9) {per-transcript\\ counts};
    \node[dry] (norm) at (2.7,0.9) {size factor\\ normalisation};
    \node[dry] (de) at (0,0.9) {negative binomial\\ Wald test};

    \foreach \a/\b in {culture/sample, sample/rna, rna/rin, rin/library,
                       library/seq} {
      \draw[flow] (\a) -- (\b);
    }
    \foreach \a/\b in {trim/align, align/mapgate, mapgate/count, count/norm,
                       norm/de} {
      \draw[flow] (\a) -- (\b);
    }
    \draw[flow] (seq.south) -- ++(0,-0.75) -- ++(1.5,0) -- ++(0,-0.9)
      -- ++(-1.5,0) -- (trim.north);

    \draw[flow,draw=accent,dash pattern=on 2pt off 1.5pt]
      (rin.north) -- ++(0,0.65) -| node[pos=0.25,above,text=accent]
      {re-extract on failure} (rna.north);
    \draw[flow,draw=accent,dash pattern=on 2pt off 1.5pt]
      (mapgate.south) -- ++(0,-0.65) -| node[pos=0.25,below,text=accent]
      {inspect unmapped fraction} (align.south);

    \node[dry] (report) at (0,-1.3) {ranked transcript list\\ and enrichment};
    \draw[flow] (de) -- (report);

    \begin{scope}[on background layer]
      \node[draw=wetlab!45,dashed,rounded corners=3pt,inner sep=4mm,
        fit=(culture)(seq)(rin)] (wetbox) {};
      \node[draw=drylab!45,dashed,rounded corners=3pt,inner sep=4mm,
        fit=(trim)(de)(mapgate)(report)] (drybox) {};
    \end{scope}
    \node[lane,text=wetlab,anchor=west] at ($(wetbox.north west)+(1mm,2.5mm)$)
      {Bench};
    \node[lane,text=drylab,anchor=west] at ($(drybox.north west)+(1mm,2.5mm)$)
      {Analysis};
  \end{tikzpicture}}
  \caption{Experimental and computational workflow. The two red boxes are the
  acceptance checks written into the protocol, and the dashed paths are the
  actions taken when a check fails.}
  \label{fig:workflow}
\end{figure}

The commands below reproduce the analysis from the raw archive. Tool versions
are pinned in the container manifest distributed with this report.

\begin{lstlisting}[style=shellstyle,caption={Pipeline invocation for one
library. The loop over the remaining five differs only in the sample
identifier.},label={lst:pipeline},captionpos=b]
fastp -i raw/NR-1_R1.fastq.gz -I raw/NR-1_R2.fastq.gz \
      -o trim/NR-1_R1.fastq.gz -O trim/NR-1_R2.fastq.gz \
      --detect_adapter_for_pe --cut_tail --cut_mean_quality 20 \
      --length_required 50 --json qc/NR-1.fastp.json

STAR --genomeDir index/tpseudo_v3 --readFilesCommand zcat \
     --readFilesIn trim/NR-1_R1.fastq.gz trim/NR-1_R2.fastq.gz \
     --outSAMtype BAM SortedByCoordinate \
     --outFilterMismatchNoverLmax 0.04 --quantMode GeneCounts \
     --outFileNamePrefix bam/NR-1.

samtools index bam/NR-1.Aligned.sortedByCoord.out.bam
featureCounts -p -s 2 -a ref/tpseudo_v3.gtf -o counts/NR-1.txt \
              bam/NR-1.Aligned.sortedByCoord.out.bam

multiqc qc bam counts -o qc/report
Rscript scripts/deseq_contrast.R --design design.csv \
        --contrast regime,limited,replete --alpha 0.05
\end{lstlisting}

\subsection{Statistical model}

Counts $K_{gi}$ for transcript $g$ in library $i$ were modelled as negative
binomial with a library-specific size factor $s_i$ and a gene-specific
dispersion $\alpha_g$,
\begin{equation}
  \label{eq:nb}
  K_{gi} \sim \mathrm{NB}\!\left(\mu_{gi}, \alpha_g\right),
  \qquad
  \mu_{gi} = s_i \, q_{gi},
  \qquad
  \log_{2} q_{gi} = \beta_{g0} + \beta_{g1} x_i,
\end{equation}
where $x_i$ indicates the nitrogen-limited regime \citep{love2014deseq2}.
Size factors were the median ratio to the geometric mean across libraries.
Dispersions were shrunk toward a fitted mean and dispersion trend, and the
regime coefficient was tested with the Wald statistic
\begin{equation}
  \label{eq:wald}
  W_g = \frac{\hat{\beta}_{g1}}{\mathrm{SE}(\hat{\beta}_{g1})},
\end{equation}
with $p$ values adjusted across transcripts by the Benjamini and Hochberg
procedure \citep{benjamini1995controlling}. Abundance is reported as transcripts
per million,
\begin{equation}
  \label{eq:tpm}
  \mathrm{TPM}_{gi}
  = 10^{6} \cdot
    \frac{K_{gi} / \ell_{g}}{\sum_{h} K_{hi} / \ell_{h}},
\end{equation}
with $\ell_g$ the effective length of transcript $g$. A transcript was called
differentially expressed when the adjusted $p$ value fell below 0.05 and the
absolute shrunken $\logfc$ exceeded 1.0.
